\documentclass[11pt]{article}

\usepackage[a4paper,margin=1in]{geometry}
\usepackage{amsmath,amssymb,amsthm}
\usepackage{mathtools}
\usepackage{hyperref}

\title{\bf Second-Order Perturbative Expansion of a Layered Clifford Circuit with Small Unitary Insertions}
\author{Notes (self-contained)}
\date{\today}

\newcommand{\id}{\mathbb{I}}
\newcommand{\cE}{\mathcal{E}}

\begin{document}
\maketitle

\section{Problem setup}

Consider an $n$-qubit layered circuit of depth $L$ whose ideal evolution is a product of layer unitaries
\begin{equation}
U \;=\; \prod_{l=1}^{L} u_l,
\label{eq:idealU}
\end{equation}
where each $u_l$ is a Clifford unitary (the Clifford assumption is only needed for later applications; the algebra below holds for general unitaries). We use the convention that the product is time-ordered from early to late: the factor with smaller $l$ is earlier in time.

We perturb each layer by inserting a small unitary generated by a Hermitian operator $H_l$:
\begin{equation}
u_l' \;=\; e^{-i\epsilon_l H_l}\, u_l,
\qquad |\epsilon_l|\ll 1,
\qquad H_l = H_l^\dagger.
\label{eq:layerPert}
\end{equation}
The perturbed circuit is therefore
\begin{equation}
U' \;=\; \prod_{l=1}^{L} u_l'
\;=\; \prod_{l=1}^{L}\Big(e^{-i\epsilon_l H_l}u_l\Big).
\label{eq:UpDef}
\end{equation}

\paragraph{Goal.}
Derive (i) a second-order expansion of $U'$ in $\{\epsilon_l\}$, and (ii) a second-order expansion of the corresponding unitary channel
\begin{equation}
\cE_L(\rho) \;=\; U'\rho\, U'^\dagger,
\label{eq:channelDef}
\end{equation}
in a form that a reader can replicate.

\section{Notation for partial products}

Define partial products of the ideal layer unitaries:
\begin{equation}
U_{a:b} \;:=\; \prod_{l=a}^{b} u_l,\qquad (a\le b),
\label{eq:Uab}
\end{equation}
and adopt the convention that an empty product equals the identity:
\begin{equation}
U_{a:b}=\id \quad \text{if } a>b..
\label{eq:empty}
\end{equation}
Then $U = U_{1:L}$. For each layer $l$,
\begin{equation}
U_{l+1:L} = \prod_{j=l+1}^{L}u_j,\qquad
U_{1:l} = \prod_{j=1}^{l}u_j.
\end{equation}

\section{Pushing insertions to the circuit end}

Conjugate each insertion $e^{-i\epsilon_l H_l}$ through the \emph{later} ideal evolution. Define the ``Heisenberg-pushed'' generators
\begin{equation}
\widetilde H_l \;:=\; U_{l+1:L}\, H_l \, U_{l+1:L}^\dagger.
\label{eq:Htilde}
\end{equation}
Then one can rewrite the perturbed circuit \eqref{eq:UpDef} as
\begin{equation}
\boxed{
U' \;=\; \Big(\prod_{l=1}^{L} e^{-i\epsilon_l \widetilde H_l}\Big)\, U.
}
\label{eq:UpPush}
\end{equation}

\begin{proof}[Derivation of \eqref{eq:UpPush}]
Use the identity $e^{-i\epsilon H}W = W\,e^{-i\epsilon W^\dagger H W}$ for any unitary $W$, applied iteratively with $W=U_{l+1:L}$. This replaces each $H_l$ by $\widetilde H_l$ as in \eqref{eq:Htilde}, and factors out the ideal product $U=U_{1:L}$ on the right, giving \eqref{eq:UpPush}.
\end{proof}

\paragraph{Remark (Clifford specialization).}
If each $u_l$ is Clifford and each $H_l$ is a Pauli (or low-weight sum of Paulis), then $\widetilde H_l$ is again a Pauli (or the conjugated low-weight sum).

\section{Second-order expansion of the perturbed unitary}

Define
\[
K \;:=\; \prod_{l=1}^{L} e^{-i\epsilon_l \widetilde H_l},
\qquad \text{so that } \; U' = K U.
\]
Using
\begin{equation}
e^{-i\epsilon_l \widetilde H_l}
= \id - i\epsilon_l \widetilde H_l - \frac{\epsilon_l^2}{2}\widetilde H_l^2 + O(\epsilon_l^3),
\label{eq:expSingle}
\end{equation}
and keeping terms up to total degree $2$ in the set $\{\epsilon_l\}$, we obtain
\begin{equation}
\boxed{
U' \;=\; \Big[
\id
- i\sum_{l=1}^{L}\epsilon_l \widetilde H_l
-\frac12\sum_{l=1}^{L}\epsilon_l^2 \widetilde H_l^2
-\sum_{1\le i<j\le L}\epsilon_j\epsilon_i\,\widetilde H_j\widetilde H_i
\Big]\,U \;+\; O(\epsilon^3).
}
\label{eq:Up2nd}
\end{equation}
Here $O(\epsilon^3)$ denotes terms of total degree at least $3$ in the $\epsilon_l$'s. The cross-term is time-ordered: $j>i$ implies $\widetilde H_j$ appears to the left of $\widetilde H_i$.

\section{Second-order expansion of the unitary channel}

Define the ideal-output state
\begin{equation}
\rho_0 \;:=\; U\rho U^\dagger.
\label{eq:rho0}
\end{equation}
Then the perturbed channel \eqref{eq:channelDef} becomes
\begin{equation}
\cE_L(\rho) = U'\rho U'^\dagger = K\rho_0 K^\dagger.
\label{eq:channelK}
\end{equation}

Write a degree expansion of $K$:
\begin{equation}
K = \id + A + B + O(\epsilon^3),
\label{eq:KAB}
\end{equation}
where the first- and second-order parts are
\begin{align}
A &= -i\sum_{l=1}^{L}\epsilon_l \widetilde H_l,
\label{eq:Adef}\\
B &= -\frac12\sum_{l=1}^{L}\epsilon_l^2 \widetilde H_l^2
-\sum_{1\le i<j\le L}\epsilon_j\epsilon_i\,\widetilde H_j\widetilde H_i.
\label{eq:Bdef}
\end{align}
Since each $\widetilde H_l$ is Hermitian, $A^\dagger=-A$.

Substitute \eqref{eq:KAB} into \eqref{eq:channelK} and keep terms up to $O(\epsilon^2)$:
\begin{equation}
K\rho_0 K^\dagger
=
\rho_0
+ A\rho_0 + \rho_0A^\dagger
+ B\rho_0 + \rho_0B^\dagger
+ A\rho_0 A^\dagger
+ O(\epsilon^3).
\label{eq:channelExpand}
\end{equation}

\subsection{First order}

Using $A^\dagger=-A$,
\begin{equation}
\boxed{
\cE_L(\rho) = \rho_0 - i\sum_{l=1}^{L}\epsilon_l\,[\widetilde H_l,\rho_0] + O(\epsilon^2).
}
\label{eq:channel1st}
\end{equation}

\subsection{Second order}

Separate the $O(\epsilon^2)$ correction into (i) ``self'' terms $\propto \epsilon_l^2$ and (ii) cross terms $\propto \epsilon_i\epsilon_j$ with $i<j$.

\paragraph{Self terms.}
Collect the contributions proportional to $\epsilon_l^2$ from $B\rho_0+\rho_0B^\dagger + A\rho_0A^\dagger$ and simplify:
\begin{equation}
\boxed{
\cE_L(\rho)\Big|_{\text{self},\,O(\epsilon^2)}
= -\frac12\sum_{l=1}^{L}\epsilon_l^2 \,[\widetilde H_l,[\widetilde H_l,\rho_0]].
}
\label{eq:self2nd}
\end{equation}

\paragraph{Cross terms.}
For $i<j$, the second-order cross contribution can be written in an anticommutator-commutator form:
\begin{equation}
\boxed{
\cE_L(\rho)\Big|_{\text{cross},\,O(\epsilon^2)}
= -\sum_{1\le i<j\le L}\epsilon_i\epsilon_j\;\{\widetilde H_j,\,[\widetilde H_i,\rho_0]\},
}
\label{eq:cross2nd}
\end{equation}
where $\{A,B\}=AB+BA$.

\paragraph{Full second-order channel.}
Combining \eqref{eq:channel1st}, \eqref{eq:self2nd}, and \eqref{eq:cross2nd} yields
\begin{equation}
\boxed{
\cE_L(\rho)
= \rho_0
- i\sum_{l=1}^{L}\epsilon_l\,[\widetilde H_l,\rho_0]
-\frac12\sum_{l=1}^{L}\epsilon_l^2 \,[\widetilde H_l,[\widetilde H_l,\rho_0]]
-\sum_{1\le i<j\le L}\epsilon_i\epsilon_j\;\{\widetilde H_j,\,[\widetilde H_i,\rho_0]\}
+ O(\epsilon^3).
}
\label{eq:channel2nd}
\end{equation}

\section{Replication checklist}

A reader can replicate the derivation via:
\begin{enumerate}
\item Fix the product convention and define $U_{a:b}$ with the empty-product rule \eqref{eq:empty}.
\item Rewrite $U'$ as $U'=KU$ by conjugating each insertion through the later ideal evolution, giving $\widetilde H_l$ in \eqref{eq:Htilde}.
\item Expand each factor $e^{-i\epsilon_l\widetilde H_l}$ to second order and multiply, keeping only terms of total degree $\le 2$, to obtain \eqref{eq:Up2nd}.
\item For the channel, set $\rho_0=U\rho U^\dagger$ and expand $K\rho_0K^\dagger$ using \eqref{eq:channelExpand}; simplify using commutator/anticommutator identities to obtain \eqref{eq:channel2nd}.
\end{enumerate}



\section{Second-Order Weak-Noise Expansion for Correlated Classical Pauli Hamiltonian Noise}

We derive a \emph{replicable} second-order (weak-noise) expression for the \emph{noise-averaged} output density matrix of a layered $n$-qubit circuit driven by ideal layer unitaries (typically Clifford), in the presence of \emph{correlated classical stochastic Hamiltonian noise} consisting of single-qubit dephasing terms and pairwise correlated dephasing terms.

\paragraph{Output.}
Given an input state $\rho$, the ideal circuit output is $\rho_0=U\rho U^\dagger$. We will show that the \emph{ensemble-averaged} noisy output takes the universal second-order form
\begin{equation}
\boxed{
\langle\rho_{\rm out}\rangle
=
\rho_0
- i\,[H_{\rm av},\rho_0]
-\frac12\sum_{a,b}\Gamma_{ab}\,[P_a,[P_b,\rho_0]]
+ O(\text{noise}^3).
}
\label{eq:main}
\end{equation}
Here $\{P_a\}$ is a Pauli-string basis (restricted to the support actually generated by the toggling frame), $H_{\rm av}$ is a (possibly zero) mean drift Hamiltonian, and $\Gamma_{ab}$ is a coefficient matrix given by double time integrals of the \emph{toggling-frame} noise correlation functions. All steps are spelled out below.

\section{Model}

\subsection{Ideal layered circuit}
Let the ideal circuit consist of $L$ instantaneous layer unitaries
\begin{equation}
U=\prod_{l=1}^{L}u_l,
\end{equation}
with total duration $T=L\tau$, where each layer occupies an interval $I_l=[(l-1)\tau,l\tau)$.
(Only the piecewise-constant nature matters; ``Clifford'' is not required for the algebra, but becomes useful because it maps Paulis to Paulis under conjugation.)

Let $U_0(t)$ denote the ideal (noise-free) evolution up to time $t$, i.e.,
\begin{equation}
U_0(t)=u_{l(t)}u_{l(t)-1}\cdots u_1\quad\text{for } t\in I_{l(t)}.
\end{equation}
Thus $U_0(T)=U$.

\subsection{Correlated classical Hamiltonian noise}
During layer $l$ (i.e., for $t\in I_l$), the system experiences an additive stochastic Hamiltonian
\begin{equation}
H_e(t)=\sum_{i=1}^n \beta_i(t)\,Z_i \;+\; \sum_{1\le i<j\le n} J_{ij}(t)\,Z_iZ_j.
\label{eq:He}
\end{equation}
The real-valued processes $\beta_i(t)$ and $J_{ij}(t)$ may be \emph{temporally correlated} and \emph{spatially correlated}:
\begin{align}
\langle \beta_i(t)\beta_{i'}(t')\rangle - \langle \beta_i(t)\rangle\langle \beta_{i'}(t')\rangle &\neq 0,\\
\langle J_{ij}(t)J_{i'j'}(t')\rangle - \langle J_{ij}(t)\rangle\langle J_{i'j'}(t')\rangle &\neq 0,
\end{align}
and similarly for cross-correlations $\langle \beta_i(t)J_{i'j'}(t')\rangle$. Angle brackets $\langle\cdot\rangle$ denote ensemble average.

\paragraph{Weak-noise assumption.}
We assume amplitudes are small enough that a second-order (Dyson) truncation is accurate:
\begin{equation}
\int_0^T \|H_e(t)\|\,dt \ll 1.
\end{equation}

\subsection{Noisy evolution and the averaged channel}
For a given noise realization, the lab-frame time-ordered propagator is
\begin{equation}
U_{\rm tot}(T)=\mathcal{T}\exp\Big(-i\int_0^T \big(H_{\rm ctrl}(t)+H_e(t)\big)\,dt\Big),
\end{equation}
where $H_{\rm ctrl}(t)$ generates the ideal layer unitaries. We will work in the toggling frame of the ideal evolution, so that only $H_e$ appears.

The corresponding (random) unitary channel is $\rho\mapsto U_{\rm tot}(T)\rho U_{\rm tot}^\dagger(T)$. The \emph{physical} noisy channel of interest is the ensemble average
\begin{equation}
\langle \cE(\rho)\rangle := \big\langle U_{\rm tot}(T)\rho U_{\rm tot}^\dagger(T)\big\rangle.
\end{equation}

\section{Toggling frame reduction}

Define the toggling-frame (interaction-picture) noise Hamiltonian
\begin{equation}
\widetilde H_e(t):=U_0^\dagger(t)\,H_e(t)\,U_0(t).
\label{eq:Ht}
\end{equation}
Then the total propagator factorizes as
\begin{equation}
U_{\rm tot}(T)=U\,U_I(T),
\qquad
U_I(T)=\mathcal{T}\exp\Big(-i\int_0^T \widetilde H_e(t)\,dt\Big).
\label{eq:UI}
\end{equation}
Hence for any input $\rho$, letting $\rho_0:=U\rho U^\dagger$,
\begin{equation}
\rho_{\rm out}
=
U_{\rm tot}(T)\rho U_{\rm tot}^\dagger(T)
=
U_I(T)\,\rho_0\,U_I^\dagger(T).
\label{eq:rhoout_toggle}
\end{equation}
Therefore it suffices to expand the toggling-frame channel generated by $\widetilde H_e(t)$ acting on $\rho_0$.

\paragraph{Pauli expansion (Clifford case).}
If $U_0(t)$ is Clifford for all $t$ (instantaneous Clifford layers), then for each $t$ the conjugates
$U_0^\dagger(t)Z_iU_0(t)$ and $U_0^\dagger(t)Z_iZ_jU_0(t)$ are Pauli strings (possibly $\pm X,\pm Y,\pm Z$ on each qubit).
Thus $\widetilde H_e(t)$ can be written as
\begin{equation}
\widetilde H_e(t)=\sum_a \xi_a(t)\,P_a,
\label{eq:Ht_pauli}
\end{equation}
for a (finite) set of Pauli strings $\{P_a\}$ and real stochastic coefficients $\xi_a(t)$ obtained by applying the Clifford ``Pauli frame'' to the original processes $\beta_i,J_{ij}$.

\section{Dyson expansion to second order}

Let
\begin{equation}
A := \int_0^T \widetilde H_e(t)\,dt,
\qquad
B := \int_0^T\!dt_2\int_0^{t_2}\!dt_1\; \widetilde H_e(t_2)\widetilde H_e(t_1).
\end{equation}
The time-ordered exponential admits the Dyson expansion
\begin{equation}
U_I(T)=\id - iA - B + O(\text{noise}^3),
\qquad
U_I^\dagger(T)=\id + iA - B^\dagger + O(\text{noise}^3),
\label{eq:dyson}
\end{equation}
where $B^\dagger=\int_0^T dt_2\int_0^{t_2}dt_1\,\widetilde H_e(t_1)\widetilde H_e(t_2)$ because $\widetilde H_e(t)$ is Hermitian.

Insert \eqref{eq:dyson} into \eqref{eq:rhoout_toggle} and keep terms through second order:
\begin{align}
\rho_{\rm out}
&= (\id - iA - B)\,\rho_0\,(\id + iA - B^\dagger) + O(\text{noise}^3)\nonumber\\
&= \rho_0 - i[A,\rho_0] + \big(A\rho_0A - B\rho_0 - \rho_0B^\dagger\big) + O(\text{noise}^3).
\label{eq:rhoout2}
\end{align}

\section{Ensemble average and the universal second-order form}

Take expectation of \eqref{eq:rhoout2}. Define the mean toggling-frame Hamiltonian
\begin{equation}
\widetilde H_{\rm av}(t):=\langle \widetilde H_e(t)\rangle,
\qquad
H_{\rm av}:=\int_0^T \widetilde H_{\rm av}(t)\,dt=\langle A\rangle.
\label{eq:Hav}
\end{equation}
Also define fluctuations $\Delta \widetilde H_e(t):=\widetilde H_e(t)-\langle \widetilde H_e(t)\rangle$ and $\Delta A:=A-\langle A\rangle$.

\subsection{First order}
The averaged first-order term is
\begin{equation}
\langle -i[A,\rho_0]\rangle = -i[\langle A\rangle,\rho_0] = -i[H_{\rm av},\rho_0].
\end{equation}
If the noise is zero-mean in the toggling frame (often the case), then $H_{\rm av}=0$.

\subsection{Second order: double-commutator structure}
A standard identity shows that
\begin{equation}
A\rho_0A - \frac12\{A^2,\rho_0\} = -\frac12[A,[A,\rho_0]].
\label{eq:doublecomm_id}
\end{equation}
Using $B+B^\dagger=A^2$ (the two time-ordered domains partition the square), one finds
\begin{equation}
A\rho_0A - B\rho_0 - \rho_0B^\dagger
=
-\frac12[A,[A,\rho_0]].
\label{eq:second_compact}
\end{equation}
Therefore, for each realization,
\begin{equation}
\rho_{\rm out}
=
\rho_0 - i[A,\rho_0] - \frac12[A,[A,\rho_0]] + O(\text{noise}^3).
\label{eq:rhoout2_compact}
\end{equation}

Now average. Write $A=\langle A\rangle + \Delta A$ and keep terms through second order:
\begin{equation}
\boxed{
\langle\rho_{\rm out}\rangle
=
\rho_0
- i[H_{\rm av},\rho_0]
-\frac12\langle [\Delta A,[\Delta A,\rho_0]]\rangle
+ O(\text{noise}^3).
}
\label{eq:rhoout_avg_compact}
\end{equation}

\subsection{Pauli-basis coefficient matrix}
Assume the Pauli expansion \eqref{eq:Ht_pauli}:
\[
\widetilde H_e(t)=\sum_a \xi_a(t)P_a,\qquad P_a\ \text{Pauli strings}.
\]
Then
\begin{equation}
\Delta A = \int_0^T \Delta\widetilde H_e(t)\,dt
= \sum_a \Big(\int_0^T \Delta\xi_a(t)\,dt\Big) P_a.
\end{equation}
Plugging into \eqref{eq:rhoout_avg_compact} gives the main result \eqref{eq:main} with
\begin{equation}
\boxed{
\Gamma_{ab}
=
\int_0^T\!dt\int_0^T\!dt'\;
\Big\langle \Delta\xi_a(t)\,\Delta\xi_b(t')\Big\rangle.
}
\label{eq:Gamma}
\end{equation}

\section{How to compute $\xi_a(t)$ from the original processes}
In the model \eqref{eq:He}, each term is either $Z_i$ or $Z_iZ_j$. Under a (piecewise-constant) Clifford Pauli frame, for each $t$ there exist Pauli strings $P_i(t)$ and $P_{ij}(t)$ such that
\begin{equation}
U_0^\dagger(t) Z_i U_0(t)=P_i(t),\qquad
U_0^\dagger(t) Z_iZ_j U_0(t)=P_{ij}(t).
\end{equation}
Hence
\begin{equation}
\widetilde H_e(t)=\sum_i \beta_i(t)\,P_i(t)+\sum_{i<j}J_{ij}(t)\,P_{ij}(t).
\end{equation}
To match \eqref{eq:Ht_pauli}, fix a Pauli basis $\{P_a\}$ containing all $P_i(t),P_{ij}(t)$ that occur (over all $t$). Then $\xi_a(t)$ is obtained by mapping each occurring Pauli string to its basis label $a$ (including the $\pm$ sign), and assigning the corresponding coefficient $\beta_i(t)$ or $J_{ij}(t)$.

\section{Replication checklist}
A reader can reproduce the result by:
\begin{enumerate}
\item Defining $U_0(t)$ and the toggling-frame Hamiltonian $\widetilde H_e(t)=U_0^\dagger(t)H_e(t)U_0(t)$.
\item Writing $U_I(T)=\mathcal{T}\exp(-i\int_0^T \widetilde H_e(t)\,dt)$ and performing the Dyson expansion to second order \eqref{eq:dyson}.
\item Multiplying $(\id-iA-B)\rho_0(\id+iA-B^\dagger)$ and simplifying to obtain \eqref{eq:rhoout2}.
\item Using $B+B^\dagger=A^2$ and the identity \eqref{eq:doublecomm_id} to reach the compact form \eqref{eq:rhoout2_compact}.
\item Ensemble-averaging and expanding in a Pauli basis to obtain \eqref{eq:main} with $\Gamma_{ab}$ given by \eqref{eq:Gamma}.
\end{enumerate}

\section{Google's tech details in RL interface}

This section summarizes the RL data interface used in \emph{Reinforcement learning control of quantum error correction} (Google Quantum AI and collaborators, Dec. 8, 2025), aligned with the $(a_t,s_t,r_t)$ notation.

\subsection*{Action $a_t$}
At learning epoch $t$, the agent samples full control-policy vectors from a Gaussian distribution over all learnable control parameters. A sampled vector is applied to the classical controller and kept fixed while collecting a batch of QEC shots/cycles. Therefore $a_t$ is not a per-gate online action; it is a full parameter setting for one candidate policy.

\subsection*{Observation $s_t$}
The direct environment feedback is detector data from QEC:
\begin{itemize}
\item binary detector events (parity flips),
\item aggregated into detector-wise detection rates (DR) over finite shots/cycles.
\end{itemize}
Operationally, for candidate $k$ in epoch $t$, one can view
\[
s_{t,k}=\big(\widehat{\mathrm{DR}}_{t,k,1},\ldots,\widehat{\mathrm{DR}}_{t,k,m}\big),
\]
where $m$ is the number of detector groups (after optional time-translation grouping).

Important detail: in this implementation the policy is mainly parameter-exploring (distribution update over control vectors), not an explicitly state-conditioned neural policy of the form $a_t=\pi_\theta(s_t)$.

\subsection*{Reward $r_t$ from $s_t$}
Google optimizes a surrogate objective based on detector statistics:
\[
C=\hat{\mathbb{E}}[D],
\]
the average detection-event rate (or equivalently a multi-objective detector-wise form). Lower detector-event rates imply better reward. A consistent scalar choice is
\[
r_{t,k}=-C_{t,k},
\]
and in multi-objective form one can use $r_{t,k,i}=-\widehat{\mathrm{DR}}_{t,k,i}$ with factor-graph-aware gradient masking.

Decoded logical error rate (LER) is used periodically for evaluation/benchmarking; the main training signal is from detector events and DR.

\subsection*{Are Steane stabilizer histories sufficient to reproduce this flow?}
Yes, in principle: full stabilizer measurement histories are sufficient up to deterministic transformations.

Given history data plus schedule metadata, detector-event streams can be constructed by parity/XOR rules (e.g., time-difference events and boundary detectors), then aggregated into DR statistics used as $s_t$ and into reward via $r=-C$.

For practical equivalence to the Google flow, the following pieces are needed:
\begin{enumerate}
\item A complete detector definition for the specific circuit (including boundary/time-edge detectors).
\item Transformation from raw stabilizer measurements to detector events.
\item DR estimation over many shots/cycles per candidate policy.
\item Optional grouping of time-translation-equivalent detectors to keep reward dimension $O(d^2)$ for memory experiments.
\item Detector-to-control-parameter mapping (factor-graph sparsity) for efficient gradient estimation.
\end{enumerate}

Hence, raw stabilizer histories are the right primitive data; the missing layer is the detector construction and aggregation pipeline.
\end{document}
